\section{问题分析}

本题的本质是一个跨时间、跨区域的算力与能源耦合调度问题。任务侧包含执行区域和开工时刻等离散选择，能源侧包含新能源分配、购售电和储能功率等连续决策，而任务不可抢占、瞬时资源容量、分钟级执行时长与小时级能源结算又使两侧具有不同的计量口径。因此，建模时不能直接以逐小时汇总 GPU 需求替代任务，也不能把所有触及某一小时的任务均视为完整运行一小时。合理的处理方式是先构造每个任务的合法区域--开工候选域，再以任务是否与某小时相交的指示量审计瞬时容量，以实际重叠时长计算 GPU-hour、AI 电量及平均设施负荷，最后将该负荷传递至新能源、储能与电网模型。

\subsection{问题一分析}

问题一同时包含描述性统计、时间序列预测和末端基础调度，但三项工作使用的数据口径并不相同。统计阶段需要把任务到达记录转换为区域--任务类型--小时序列，以任务数量、GPU 需求、GPU-hour、执行时长、到达强度、峰值并发及周期相关性刻画异构需求。预测阶段应严格采用顺序划分，先利用第 $0$ 至第 $2351$ 小时拟合候选模型，再以第 $2352$ 至第 $2375$ 小时完成模型比较；模型结构确定后，应在规定的历史区间上重新拟合，并只将第 $2376$ 至第 $2399$ 小时作为最终测试区间。由于部分区域和任务类型可能出现零需求，WAPE 只有在真实需求绝对值之和为正时才具有意义，故 MAE、RMSE 与带分母门禁的 WAPE 应共同用于误差评价。

需求序列可能同时呈现日周期、周周期、短期惯性和突发扰动。为此，预测量被表示为区域与任务类型的逐小时到达 GPU 需求，模型中应包含短期滞后、日内滞后、周内滞后、正余弦周期项以及区域和任务类型效应。考虑到尖峰任务会显著影响平方损失，固定 Huber 回归可被用于降低异常峰值对参数估计的支配作用。预测输出在模型拟合后才执行非负截断，以避免截断操作改变训练目标。实际实现采用固定 Huber 滞后回归，未执行完整候选网格选模，因此最终测试误差只能说明该已实现模型在保留测试时段上的表现。

最后 $24$ 小时的调度不能由预测需求替代，而应直接读取同期实际到达任务。对于任务 $i$，其持续小时数由分钟时长换算得到，有效截止时点取任务有效截止与系统终点中的较早者；随后根据到达时刻、任务类型、单向时延和完成边界生成合法执行区域与整数开工时刻。实时推理任务的开工时刻被固定为到达时刻，弹性任务则可在候选域内延迟执行。实际算法采用确定性贪心启发式，根据等待时间、网络时延、峰值 GPU 利用率及固定词典序选择候选，并通过完整资源审计构造可行方案；契约中规划的有限冲突回溯并未实现，故若贪心构造失败，只能报告在既定方法内未能确认可行性。

资源约束需要区分瞬时口径和能量口径。只要任务在小时 $t$ 内存在正重叠，其持续占用的 GPU 数及相应瞬时 IT 功率就应被计入容量审计；任务在该小时的实际重叠时长则用于形成 GPU-hour 和 AI 电量。区域总 IT 功率由 AI 瞬时功率与不可迁移的非 AI IT 负荷相加，设施功率由同一总 IT 功率乘以区域 PUE 得到。甘特图应以任务实际开工和完成时点绘制，并同时给出峰值 GPU 利用率与时间加权平均利用率，从而避免仅以小时平均值掩盖短时容量冲突。

\subsection{问题二分析}

问题二需要以第 $0$ 至第 $2399$ 小时的实际任务重新进行全时域调度，并将第 $2400$ 至第 $2405$ 小时发生的末端任务能耗纳入成本和碳排放核算。任务侧沿用问题一的合法候选域、唯一执行、实时即时开工、截止时间、单向时延及瞬时容量约束，但不能直接沿用预测结果。AI 任务的小时电量由 GPU 需求、任务类型功率映射与真实重叠时长共同计算，区域设施平均负荷则由 AI 平均 IT 负荷和固定非 AI IT 负荷之和乘以 PUE 获得。

由于问题二不配置储能，逐时能源分配可在给定设施负荷后被确定性复算。新能源先被用于满足本地负荷；当新能源不足时，缺口由电网购电补足；当新能源富余时，可在区域外送上限内售出，其余部分被记为弃电。购电与售电不能在同一时段同时发生，超过进口上限的任务候选应被判定为不可行。系统成本按照购电支出扣除售电收入计算，碳排放只由电网购电量与对应碳强度相乘累积，新能源利用率的分子包含直接使用、允许计入的新能源充电和新能源外送，而分母为可用新能源总量。

成本、碳排放、网络时延、新能源未利用程度与等待量具有不同量纲，因而不能直接相加。方法契约给出的处理是以相对基准真实差值的有限正中位数作为尺度，对有效维度进行标准化后加权评价；若某一指标不存在正尺度或其定义分母为零，该维度应退出评分。网络时延采用 GPU 需求与执行时长加权的单向时延，只有权重和为正时才计算。实际实现并未执行规定数量的单任务迁移或开工邻域改进，而是对确定性任务方案进行完整能源复算，因此所得方案属于启发式可行方案，不能被解释为连续参数或离散候选空间中的最优解。

\subsection{问题三分析}

问题三与前两问的关键区别在于任务负荷已经固定。区域 IT 负荷应直接取基准 AI IT 负荷与非 AI IT 负荷之和，设施负荷再乘以区域 PUE，任务迁移和开工时刻不得被重新优化。这样，问题三可被转化为逐区域、逐小时的能源与储能协同问题。无储能状态只用于建立比较基准，即使该状态违反购电边界，也不能据此截断含储能方案的搜索；含储能方案必须从附件给定的绝对初始 SOC 独立向前构造。

能源侧需要同时满足新能源守恒、负荷平衡和储能动态。可用新能源被拆分为直接使用、储能充电、外送和弃电，电网购电则被拆分为直接供负荷和储能充电。储能总充电功率等于新能源充电与电网充电之和，时段末 SOC 由前一状态、充电效率和放电效率递推。每次动作后还需检查未来剩余最大充电能力能否保证终端 SOC 不低于初始 SOC，以避免前期过度放电造成末端不可恢复。充电、放电和空闲三种状态相互排斥，购电与售电也应通过状态变量或等价逻辑保持互斥。

含储能与无储能状态的比较需要统一指标口径。净购电被定义为购电减售电，区域峰值指标取整个时域净购电的非负最大值，负荷波动采用相邻小时净购电绝对变化量之和，等效循环次数仅在储能容量为正时计算。运行成本、碳排放和新能源利用率均应由各自方案的能源流重新计算。实际实现采用规则型 SOC 前向策略，并未执行完整动作邻域搜索，因而储能策略只能被表述为满足已审计边界的规则型候选；若含储能成本高于无储能状态，也不能将其包装为经济性改进。

\subsection{问题四分析}

问题四把前三问中的任务调度、能源平衡与储能动态统一起来。任务区域和开工时刻决定 AI 小时电量及设施负荷，设施负荷进一步决定新能源消纳、电网购电与储能动作，能源价格和碳强度又会反向影响弹性任务的迁移和延迟选择。由于离散任务决策与连续能源决策跨越完整时域，直接穷举全部组合并不可行，因此应先删除违反时延、到达、截止和资源下界的候选，再独立构造全时域联合可行初始方案，并通过有界交替邻域对任务与储能动作进行改进。

联合评价包含运行成本、购电碳排放、有效网络时延、有效等待损失、新能源未利用率及六区域非负峰值净购电功率。服务质量中的实时即时开工率、SLA 满足率和按期完成率本质上属于硬约束审计，只有分母为正且候选之间存在差异的等待分量才适合作为评分项。实际实现以平均等待构造有限代理，并未实现完整等待质量公式。各指标应采用有限正差值尺度进行无量纲化，分母无效或候选间无变化的指标不得被强行纳入综合评分。

情景比较需要保持单因素可比。不同碳约束、电价机制和新能源水平或波动情景均应从同一联合基准独立重建，在相同任务集、容量、时延、PUE 和功率映射下复算全部指标。随机新能源波动使用固定种子，以保证重复运行的一致性。实际方法未执行约定的共享交替邻域搜索，每个情景只重建并审计一个规则型候选，且问题四的任务迁移方案保持联合基准不变。因此，未达到某一碳目标只能说明声明的启发式候选未满足该目标，不能据此推断该情景在数学上不可行；情景结果也只构成规则型策略比较，不构成 Pareto 前沿。

\subsection{子问题间的衔接关系}

四个子问题共享任务重叠计算、功率换算与约束审计口径，但其决策变量和输入负荷不能相互混用。问题一的预测模型可独立训练，预测输出只服务于精度评价，最后 $24$ 小时基础调度仍由实际任务生成。问题二可以复用问题一的候选域与容量审计方法，但必须以全时域实际任务和逐时能源数据重新计算。问题三固定采用附件中的基准 AI IT 负荷，因而可与任务预测和迁移调度相对独立地求解。问题四则继承任务约束、能源守恒、储能状态和统一指标定义，但联合方案必须独立构造，不能简单拼接问题二的任务方案与问题三的储能方案。
